\documentclass[12pt]{article}
\usepackage{setspace}
\usepackage{amsmath,amssymb,euscript,graphicx}
\usepackage{xurl}
\usepackage{graphicx}
\usepackage{subfigure}
\usepackage{wrapfig}
\usepackage[compact]{titlesec}
\usepackage{fancybox,color}
\usepackage[footnotesize]{caption}
\usepackage{cite}
\usepackage{enumitem}
\usepackage[normalem]{ulem}
\usepackage[top=0.9375in, right=0.9375in, left=0.9375in, bottom=0.9375in, centering]{geometry}
\usepackage{bm}

\usepackage{amsthm}


% -- Loading the code block package:
\usepackage{listings}
% -- Basic formatting
\usepackage[utf8]{inputenc}
\usepackage[english]{babel}
\usepackage{times}
\setlength{\parindent}{8pt}
\usepackage{indentfirst}
% -- Defining colors:
\usepackage[dvipsnames]{xcolor}
\definecolor{codegreen}{rgb}{0,0.6,0}
\definecolor{codegray}{rgb}{0.5,0.5,0.5}
\definecolor{codepurple}{rgb}{0.58,0,0.82}
\definecolor{backcolour}{rgb}{0.95,0.95,0.92}
% Definig a custom style:
\lstdefinestyle{mystyle}{
    backgroundcolor=\color{backcolour},   
    commentstyle=\color{codepurple},
    keywordstyle=\color{NavyBlue},
    numberstyle=\tiny\color{codegray},
    stringstyle=\color{codepurple},
    basicstyle=\ttfamily\footnotesize\bfseries,
    breakatwhitespace=false,         
    breaklines=true,                 
    captionpos=t,                    
    keepspaces=true,                 
    numbers=left,                    
    numbersep=5pt,                  
    showspaces=false,                
    showstringspaces=false,
    showtabs=false,                  
    tabsize=2
}
% -- Setting up the custom style:
\lstset{style=mystyle}


% IEEE uses Times Roman font, so we'll default to Times.
% These three commands make up the entire times.sty package.
\renewcommand{\sfdefault}{phv}
\renewcommand{\rmdefault}{ptm}
\renewcommand{\ttdefault}{pcr}

\newcommand{\squeeze}{\vspace{-0.1in}}
\newcommand{\changed}[1]{\textcolor{blue}{#1}} %generates text in blue

%%%%%%%%%%%%%%%%%%%%%%%%%%%%%%%%%%%%%%%%%%%%%%%%%%%%%%%%%%%%%%%%%%%%%%%

\begin{document}

\begin{center}
  {\Large {\bf Assignment 1: Shapes and Functions}}
  
\end{center}

\vspace{5pt}
\hrule
\vspace{5pt}

\section{Summary}
In this assignment, you will make more shapes and write functions to introduce generality to your shapes. 

\section{Creating a new shapes}

\subsection{Inventing a shape}

\subsection{Build a Python program to draw the new shape}

\subsection{ Create another new shape}

\section{Make a pattern using two shapes}

\subsection{Make functions for the two shapes}

\subsection{Create a new shape or pattern out of \texttt{shapeA} and \texttt{shapeB}}4

\section{Use parameters to control how your shapes are drawn}

\section{Extensions}

\section{Acknowledgments}
These instructions are from Professor Hannen Wolfe in the Colby College CS Department. 
Sources:
\begin{itemize}
\item \footnotesize{\url{https://cs.colby.edu/courses/F25/cs151a/projects/01/Project01WhatIsAShape.html}
\end{itemize}


\end{document}
